\section{Economic Validation of the Linkages}
\label{sec:policy}

The link is economically valid if the information passed between the models represents
the same real-world object on both sides and produces the response implied by economic
theory. This requires more than showing that a number can be transferred. The receiving
model must place it on the correct households, industries, or government account; the
direction of the response must make sense; and the same cost, investment, or revenue must
not enter through two routes.

\subsection{The economic tests}

The electricity-cost relationship passes a system-cost change to the two agents that pay
for electricity: households through final consumption and firms through their direct and
indirect electricity exposure. A higher cost should reduce household purchasing power,
raise costs most in electricity-intensive industries, and impose a larger proportional
burden on households for whom electricity is a necessity. Returning the notional revenue
to households isolates this price effect; otherwise the relationship would silently add a
tax and transfer policy that did not originate in CLEWS.

The two capital relationships answer different questions. Grid investment adds public
infrastructure and requires an explicit financing assumption. Generation capital
intensity changes the way the electricity industry produces. An investment incentive,
by contrast, changes the industry's cost of acquiring capital. These effects should not
be collapsed into a single ``investment'' adjustment: a more capital-intensive technology
mix does not by itself determine how much private capital the economy will ultimately
employ.

The health relationship follows a separate causal path. Lower particulate emissions reduce
pollution-related mortality and morbidity; this changes survival, population, and effective
labour. Its economic logic is well established, but its magnitude depends on the
emissions-to-exposure and exposure-to-health translations. Validation therefore requires
both the expected demographic and labour responses and epidemiological calibration suited
to the country application.

The carbon-price relationship must distinguish a resource cost from a government payment.
A carbon price can change the CLEWS technology plan, create a payment by households or
firms, and generate public revenue. These are different consequences of the same policy.
They may be represented together, but each monetary flow must appear only once and any
revenue returned to households must have been collected first.

The return relationships are also economically interpretable, with different degrees of
directness. Economic activity is the appropriate basis for revising demand for energy
services: industrial demand should follow the activity of the industries using electricity,
while household demand should follow household consumption. Using the equilibrium return
as a planning discount rate is a harmonisation of the opportunity cost of capital across
the models. It is not an assertion that a market return, a social discount rate, and a
project financing rate are identical; taxes, risk, maturity, and policy objectives determine
the adjustment required for a particular application.

The validation runs produced the responses summarised in Table~\ref{tab:validation}.
They are evidence about the operation of the economic mechanisms, not estimates of the
benefits or costs of a particular policy.

\begin{table}[H]
\centering\scriptsize
\caption{Observed economic behaviour in the OG--CLEWS validation runs.}
\label{tab:validation}
\begin{tabularx}{\linewidth}{@{}>{\bfseries\raggedright\arraybackslash}p{0.20\linewidth} >{\raggedright\arraybackslash}p{0.27\linewidth} L@{}}
\toprule
Relationship & Economic proposition & Observed result \\
\midrule
Electricity cost & system costs affect household purchasing power and industry costs & a roughly $5\%$ cost increase reduced steady-state household consumption by $0.062\%$; with the payment recycled, the proportional burden was larger for lower-income households \\
Public grid investment & infrastructure raises public capital and must be financed & the source scenario contained almost no additional grid investment, and the public-capital and aggregate responses were correspondingly close to zero \\
Generation capital intensity & technology choice changes the electricity industry's production structure & a $12.2\%$ increase in the capital share lowered the electricity price $42.8\%$, raised sector output $20.9\%$, and reduced sector capital $22.4\%$; the firm identity reproduced the capital response within $0.03\%$ \\
Private investment incentive & policy changes the cost of private generation capital & a $20\%$ investment tax credit raised private capital in electricity by $5.0\%$, while other industries moved by less than $0.02\%$ and aggregate output was essentially unchanged \\
Health & particulate emissions affect survival and effective labour & a $2.7\%$ emissions reduction produced about $95$ lives saved and a near-term output gain of about $0.19\%$; the long-run output effect was near zero because the survival gain was concentrated among older ages \\
Carbon price & one policy changes resource costs, payments, and public revenue & the recycled standalone experiment reduced output by $0.021\%$, capital by $0.019\%$, and labour by $0.045\%$, while consumption increased by $0.026\%$ \\
Economic activity & industry and household activity determine service demand & a controlled $10\%$ increase in household electricity demand raised CLEWS electricity production by $4.28\%$ \\
Equilibrium return & the economy's opportunity cost of capital informs resource planning & the equilibrium-return path was produced as the planning-rate input; its effect on resource choices has not yet been measured in the recorded validation runs \\
\bottomrule
\end{tabularx}
\end{table}

\subsection{What the validation supports}

The observed responses follow the economic mechanisms represented by the linkages. Costs
reduce purchasing power, distribution reflects expenditure patterns, a change in factor
shares is distinct from an investment incentive, health effects differ by horizon, and
recycling changes who ultimately bears a policy cost. The results also identify where
interpretation depends on calibration. The
electricity-cost relationship uses an average-cost proxy rather than an observed tariff or
marginal price. The health effect requires a country-specific pollution and disease
translation. Monetary carbon and investment flows require a common unit and deflator.
The discount-rate relationship requires a stated choice between social, market, and
financing concepts.

These qualifications do not weaken the economic case for the linkages. They define the
evidence needed to use each one credibly. The central validation claim is therefore that
the linked system represents distinct and defensible pathways from resource decisions to
households, production, health, demography, and public finance---and back from economic
conditions to resource demand and planning assumptions.
